\section{主要功能}

\subsection{系统核心功能}

\subsubsection{单条评价分析}

用户可以输入中文、英文或中英混合评价，也可以选择系统内置示例。
页面返回语言识别结果、情感倾向、置信度、模型来源、课程维度、命中证据和关键词。
当启用建议生成时，系统会将结构化分析结果发送给 ChatECNU；
接口不可用时，自动使用本地模板生成同样结构的总结与建议。

\subsubsection{批量 CSV 分析}

系统支持上传包含 \code{review_text}、\code{text}、\code{review}、\code{comment} 或
\code{content} 等常见文本字段的 CSV 文件。
批量分析输出有效评价数量、积极反馈占比、待关注评论数量、主要关注维度、情感分布、维度排行、
高频问题关键词、典型评论证据和逐条分析明细。
分析结果可下载为 CSV，适合期末集中评价或问卷开放题的快速整理。

\subsubsection{固定案例验证}

固定案例验证页面读取 \code{data/test_cases.csv} 中的 12 条固定案例；
一键执行完整 NLP 流程，并同时比较情感标签和课程维度是否符合预期。
页面给出固定案例总数、通过案例、失败案例、通过率和逐条结果。
该功能用于课堂演示和回归检查，不作为模型泛化能力指标。

\subsubsection{模型评估}

模型评估页面展示离线实验结果和当前运行状态。
页面包括最终实验主模型、测试集 Accuracy、Macro-F1、评估样本数、配置后端、实际运行后端、
BERT 权重状态、TF-IDF 模型状态和 LLM 配置状态。
同时展示传统分类器对比、BERT 混淆矩阵、分类报告和固定案例消融实验。

\subsection{功能模块划分}

\begin{table}[H]
  \centering
  \caption{系统功能模块与代码对应关系}
  \label{tab:function-modules}
  \small
  \begin{tabularx}{\textwidth}{C{2.7cm} L{3.5cm} Y}
    \toprule
    模块 & 主要文件 & 功能说明 \\
    \midrule
    前端交互 & \code{app.py} &
    提供首页概览、单条分析、批量分析、固定案例验证和模型评估五个页面 \\
    文本预处理 & \code{preprocess.py}、\code{sentiment_normalizer.py} &
    语言识别、文本清洗、双语分词、缩写展开和受限拼写纠错 \\
    情感分析 & \code{bert_sentiment.py}、\code{nlp_analyzer.py} &
    BERT 推理、规则校正、TF-IDF 回退和统一结果编排 \\
    课程维度 & \code{topic_analyzer.py} &
    识别六类课程维度，并返回命中关键词和原文证据 \\
    关键词提取 & \code{keyword_extractor.py} &
    提取单条评价关键词和批量评价中的高频问题表达 \\
    建议生成 & \code{llm_client.py} &
    调用 ChatECNU 生成结构化建议，异常时使用本地模板 \\
    训练与评估 & \code{train_model.py}、\code{train_bert.py}、\code{scripts/} &
    模型训练、指标保存、消融实验、压力测试和报告图表生成 \\
    \bottomrule
  \end{tabularx}
\end{table}

各功能模块在系统中的层次关系见图~\ref{fig:system-architecture}。
从评价输入到结构化结果和教学建议的处理顺序见图~\ref{fig:data-flow}。
模块表、总体架构图和数据流图共同说明了页面功能、后端服务与模型资源之间的对应关系。

\subsection{功能设计原则}

系统功能设计遵循“先结构化分析，再生成自然语言建议”的原则。
情感标签由可评估的分类模型与规则模块决定；
大语言模型只接收情感、课程维度、关键词和证据等结构化结果。
这样可以减少生成模型对核心标签的不可控影响，也能够在没有 API 时保持分析结果一致。

第二个原则是“输出结论必须尽量附带证据”。
课程维度不是只返回一个名称，而是同时记录命中关键词、关键词位置和附近的原文片段。
教师可以从“系统判断”快速回到“学生原话”，降低错误归纳带来的风险。
